\section{Связь $(3{+}1)$ и $(2{+}1)$}

Гексагональный слой $(2{+}1)$ --- срез пространства Минковского $(3{+}1)$ с FCC
(\cref{fig:carrier-slice-hex-layer,fig:carrier-slice-fcc-layer}).
Мост $\speedC=\caKappaGeom \caSignalSpeed$ с $\caKappaGeom=1/\sqrt{2}$ якорится на полном трёхмерном носителе;
чистый $(2{+}1)$ без объёма несёт $\caKappaGeom=\sqrt{3}/2$.

\begin{figure}[htbp]
  \centering
  \begin{tikzpicture}[carrier panel]
  \CarrierHexRing{0.38}{0.65}
  \filldraw[carrier object] (0,0) circle (2pt);
\end{tikzpicture}

  \caption{Гексагональный срез $(2{+}1)$: $\caKappaGeom_{\mathrm{hex}}=\sqrt{3}/2$, $|N|=6$.}
  \label{fig:carrier-slice-hex-layer}
\end{figure}

\begin{figure}[htbp]
  \centering
  \begin{tikzpicture}[carrier panel]
  \draw[carrier object] (0,0) circle (0.95);
  \foreach \a in {0,72,...,288} {
    \filldraw[carrier object, carrier fill] ({0.95*cos(\a)},{0.95*sin(\a)}) circle (2pt);
  }
  \filldraw[carrier object] (0,0) circle (2.5pt);
  \draw[carrier muted, dashed, line width=0.7pt] (-1.15,-0.55) -- (1.15,0.55);
  \node[anchor=south west, font=\scriptsize, carrier muted] at (-1.0,0.45) {$\{111\}$};
\end{tikzpicture}

  \caption{Носитель $(3{+}1)$ FCC и срез $\{111\}$: $\caKappaGeom=1/\sqrt{2}$, $|N|=12$.}
  \label{fig:carrier-slice-fcc-layer}
\end{figure}

